
% \clearpage % TEMP temp clearpage
\chapter
[\CO{[URO]} Urologische Erkrankungen]
{\CC{[URO]} Urologische Erkrankungen}
\ChapterIntro
{%
~~~
}
{%
\begin{description}
  \item[Maintainer:] Sebastian Gabriel
  \item[Autoren:] Diverse
  \item[Reviewer:] Standard-Reviewprozess
  \item[Version:] {\DocVersion} ({\DocVersionInfo})
  \item[SHA1:] {\DocGitCommit}
\end{description}

{\TxTeilkapitelAASS}
}

\Definition{Erkrankungen des harnableitenden Systems.}

\subsection{Exkurs: Harnkatheter}

\Layout{0}{\TxBeschreibung}
{Ein Harnkatheter ist ein harnableitendes System, bei dem ein
Schlauch (Katheter) über die Harnröhre in die Harnblase
eingeführt wird um Harn abzulassen. Mittels eines
aufblasbaren Ballons kann der Katheter in der Blase fixiert
werden, man spricht dann vom \I{Dauerkatheter}. Am
Harnkatheter kann ein \II{Harnbeutel} angeschlossen werden,
welcher den Urin sammelt und sich entleeren lässt. Der
Harnbeutel muss \EE{unter dem Patienten angebracht} sein, da
andernfalls der Urin zurückrinnen würde. Beim Umlagern
ist daher der zum Patienten führende \EE{Schlauch
abzuklemmen}. Der Katheter und das daran hängende System ist
eine mögliche \EE{Eintrittspforte für Keime}, daher ist auf einen
sauberen, kontaminationsfreien Umgang zu achten.

Wenn der Harnbeutel ausgelassen werden muss, so ist die
abgelassene Menge zu notieren, da bei manchen Patienten der
Flüssigkeitshaushalt bilanziert wird.

Ein Harnkatheter ist eine
medizinisch-pflegerische Routinemaßnahme und entsprechend oft bei Patienten anzutreffen.}
{
\begin{itemize}
  \item Katheter leitet aus Harnblase Urin ab
  \item Harnbeutel -- \E{unterhalb} d. Patienten
  \item Schlauch abklemmen beim Umlagern
  \item Wenn ablassen, Menge notieren
  \item Infektionsgefahr
\end{itemize}
}
{}
{}
{}



\begin{PictureSeriesWide}{}{Bilderserie: Harnkatheter}
\PictureSeriesRow{3}
{./images/demaw/Skriptumaegf-img52}
{Ein Harnkatheter \SourcePicture{DEMAW/Blacky}}
{./images/demaw/Skriptumaegf-img53}
{Ein Harnbeutel \SourcePicture{DEMAW/Blacky}}
{./images/demaw/Skriptumaegf-img59}
{Einführen eines Harnkathters
unter sterilen Bedingungen \SourcePicture{DEMAW/Blacky}}
{empty}{}
\end{PictureSeriesWide}


\subsection{Nierenkolik}
\label{topic:nierenkolik}
\index{Kolik!Nieren-}
\index{Nierenkolik}

\Definition{Blockade des Harnleiters durch Nierensteine
mit massiven kolikartigen Schmerzen aufgrund des Harnstaus.}

\Z{\emph{Vgl. Kolik: \ref{topic:kolik}}}

\Layout{0}{Ursachen}{%
Die Nierenkolik ist eine äußerst schmerzhafte
akute Erkrankung. Sie entsteht dadurch, dass Nierensteine
den Harnleiter blockieren und verstopfen. Dadurch wird der Harn
gestaut und der Harnleiter gedehnt. 
}{
\begin{itemize}
    \item Nierenstein(e)
\end{itemize}
}{}{}{}

\Layout{0}{\TxSymptome}{%
Die Folge sind massive,
kolikartige Schmerzen, in der Flanke bzw. im Unterbauch der
jeweiligen Seite, welcher auch in den Hoden bzw. in die
Schamlippen ausstrahlen kann. Im Gegensatz zu anderen
abdominellen Erkrankungen ist der Patient extrem unruhig und agitiert (\Speech{\enquote{Der Stein wandert, der Schmerz wandert, der
Patient wandert}}).
Oft läuft er gestikulierend und wehklagend durch die Gegend.
Mitunter ist eine entsprechende Vorerkrankung
\E{(Nierensteine)} bereits bekannt und kann erfragt werden.

\bigskip

\bigskip

}{
\begin{itemize}
    \item Rücken-/Flankenschmerzen, Ausstrahlung in
    Hoden/Schamlippen
    \item Evtl. Unterbauchschmerzen bei  Stein-/Sandabgang
    \item \Speech{\enquote{Der Stein wandert, der Schmerz wandert, der Patient
    wandert}}
    \item Blut im Harn (Hämaturie)
\end{itemize}
}{}{}{}



\MassnahmenMK{Spezielle
Maßnahmen}{Nierenkolik}{m:nierenkolik}{ Im
Vordergrund bei der Versorgung steht der rasche zügige Transport an eine Abteilung für Urologie. Wenn die Schmerzen
sehr massiv sind, muss eventuell eine Schmerztherapie durch
den Notarzt erfolgen.
}{}
% \MassnahmenNG{Spezielle
% Maßnahmen}{Nierenkolik}{\label{m:nierenkolik}}{ Im
% Vordergrund bei der Versorgung steht der rasche zügige Transport an eine Abteilung für Urologie. Wenn die Schmerzen
% sehr massiv sind, muss eventuell eine Schmerztherapie durch
% den Notarzt erfolgen.
% \begin{itemize}
%   \item Lagerung nach Wunsch des Patienten
%   \item Bei unerträglichen Schmerzen Notarzt zur
%   Schmerztherapie beiziehen
%   \item Transportentscheidung: Abt. für Urologie
% \end{itemize}
% }


\subsection{Harnwegsinfekt}
\label{topic:harnwegsinfekt}



\Layout{2}{\TxBeschreibung}
{Als Harnwegsinfekt wird allgemein eine meist bakterielle
Infektion der Harnwege bezeichnet. Besonders gefährdet sind
Frauen, infolge der kürzeren Harnröhre und der
normalen bakteriellen Besiedelung der Scheide. Die
häufigsten Erreger sind klassische Darmkeime \Z{(E. coli)}.

Harnwegsinfektionen werden häufig in
Gesundheitseinrichtungen übertragen (\enquote{Spitalsinfektion}).
Problematisch ist hierbei, dass bei derartigen Infektionen
häufig multiresistente Keime (\enquote{Spitalskeime},
gegen Antibiotika unempfindliche Bakterien,
\FullPrettyRef{topic:multiresistente-keime}) der Auslöser
sind, welche deutlich schwieriger zu behandeln sind. Ein
besondere Risiko stellen hierbei harnableitende Systeme
(Harnkatheter, -beutel) da.

Wird ein Harnwegsinfekt nicht behandelt, kann sich u.\,U.
die Infektion weiter ausbreiten und sich bis hin zur Sepsis
entwickeln (\FullPrettyRef{topic:sepsis}).

}
{}
{}
{}
{}



\Layout{0}{\TxSymptome}
{Der Patient klagt über Schmerzen beim Harnlassen, sowie
häufigen, aber nicht ergiebigen Harndrang. Es können
Scherzen im Nierenlager oder der Blasengegend bestehen.
Allgemeinsymptome wie Fieber, Müdigkeit und
\enquote{Krankheitsgefühl} sind möglich. 

}
{
\begin{itemize}
  \item Schmerzen beim urinieren, evtl. im Nierenlager,
  Blasengegend
  \item Häfuger, unergiebiger Harndrang
  \item Fieber
\end{itemize}}
{}
{}
{}


\subsection{Nierenbeckenentzündung}
\label{topic:nierenbeckenentzuendung}

\Z{\HeadingSub{\Lang{Syn.} \I{Pyelonephritis}}}

Die Entzündung des Nierenbeckens ist oft bakteriell bedingt 
und umfasst die oberen Harnwege, das Nierenbeckenkelchsystem
und u.\,U. auch das Nierengewebe und eine der häufigsten
Erkrankungen der Niere \cite{Pschy259}. Bei der akuten
Nierenbeckenentzündung leidet der Patient unter Fieber, Flankenschmerzen, häufigen, aber nicht ergiebigen
Harndrang sowie Schmerzen beim urinieren. Ein chronischer
Verlauf ist oft symptomarm.



\subsection{Akutes Harnverhalten}
\label{topic:harnverhalten}

\Definition{Abflussbehinderung des Harns aus der
Harnblase, verschiedene Ursachen sind möglich.}





\Layout{0}{Ursachen}
{%
Das akute Harnverhalten ist besonders bei Männern im \enquote{besten}
oder fortgeschrittenem Alter häufig anzutreffen. Hauptursache  dafür ist eine altersbedingte Vergrößerung der
Prostata, die die Harnröhre abschnürt. Davon abgesehen
können auch alte Narben, Medikamente oder eine Störung der
Schließmuskel der Harnblase\footnote{zum Beispiel im Rahmen
eines Querschnitt-Syndroms nach Wirbelsäulenverletzungen} zu
Harnverhalten führen.
}
{
\begin{itemize}
  \item Vergrößerung der \EE{Prostata}
  \item Lähmung (Querschnittlähmung)
  \item Schließmuskel-Krampf
  \item Verlegung der Harnröhre
  \item Medikamente
\end{itemize}
}{}{}{}

Je nach Füllungszustand der Harnblase kann das
Harnverhalten sehr schmerzhaft sein.
Die Maßnahmen beschränken sich auf einen zügigen Transport
an eine Abteilung für Urologie. Eine Notarzt-Nachforderung
ist kaum sinnvoll, da ein \Abk{NEF} oder \Abk{NAW} selten geeignetes
Material zur Harnableitung mitführt. 

\MassnahmenNG{Spezielle Maßnahmen}{Akutes
Harnverhalten}{m:harnverhalten}
{
}{}
% \MassnahmenNG{Spezielle Maßnahmen}{Akutes
% Harnverhalten}{\label{m:harnverhalten}}
% {
% \begin{itemize}
%     \item Zügiger Transport in eine Urologische Ambulanz
%     \item \Z{vorsichtig ablassen über Harnkatheter oder
%     Blasenpunktion (Spital)}
% \end{itemize}
% }




\subsection{Niereninsuffizienz und Nierenversagen}
\label{topic:niereninsuffizienz}
\label{topic:cni-therapie-dialyse}
\label{topic:dialysepatienten-gefahren}

\AuthNotes{Abkz. CNI}

Meistens kommt es im Rahmen anderer Krankheiten (Diabetes
mellitus, Hypertonie, erhöhte Blutfette \ldots) zu einer
massiven Schädigung der kleinen Nierengefäße und in weiterer
Folge zum Nierenversagen. Da die zuvor erwähnten Krankheiten
sehr häufig sind, ist in weiterer Folge die
Niereninsuffizienz ebenfalls häufig anzutreffen.

Als Therapie steht der Medizin anfänglich eine medikamentöse
Behandlung, später jedoch nur die Nierentransplantation oder
die \E{Blutwäsche} (\T{Dialyse}\index{Dialyse}) zur
Verfügung. Bei der klassischen Dialyse wird der Patient jeden
zweiten oder dritten Tag für einige Stunden an eine Maschine
angeschlossen, welche sein Blut reinigen kann. Die
\enquote{Dialysefahrten} sind eine der Hauptaufgaben des
Krankentransportdienstes. Dialyse-Patienten haben häufig
einen Dialyse-Shunt (\FullPrettyRef{topic:shunt}).

\index{Dialyse-Patient!Gefahren}
\EE{Der Ausfall der normalen Nierenfunktion hat Auswirkungen}
auf den Wasserhaushalt, den Elektrolythaushalt und den
Säure-Basen-Haushalt. Diese Funktionen der Niere sind zum
Überleben elementar (Vitalfunktionen zweiter Ordnung).
Dialysepatienten müssen daher sehr darauf achten nicht zu
viel zu trinken. Weiters können \E{Elektrolytstörungen} sogar
\E{lebensbedrohliche Notfälle} nach sich ziehen (bis hin zum
Herzstillstand).

Es ist daher wichtig auf den Dialysepatienten -- auch im
\enquote{normalem} Krankentransport -- immer ein wachsames Auge zu
haben.







% \subsection{Hodentorsion}
% 
% Verdrehung eine Hodens im Hodensack  ={\textgreater}
% 
% Abschnürung von Samenstrang und  Blutgefäßen
% Oft nach Eintritt in Pubertät, ev. {\textgreater} 20 J.
% 
% \paragraph{Symptome:}
% 
% {\mdseries
% akuter heftiger Schmerz}
% 
% {\mdseries
% Schwellung}
% 
% \paragraph{Therapie:}
% 
% hochlagern
% 
% umgehende urologisch-chirurgische  Intervention (ad URO!!!)




% \clearpage

% \section*{\protect\dsliterary{} Weiterführende Literatur}
% \begin{WidePar}\tiny
% \begin{multicols}{2}
% \Z{%
% \begin{labeling}{mmm}
%   \item[\citep{Renz-Polster2006}] \emph{\bibentry{Renz-Polster2006}} 
%   \item[\citep{SiewertChirurgie8}] \emph{\bibentry{SiewertChirurgie8}}
%   \item[\citep{ScholzNotfallmedizin2}] \emph{\bibentry{ScholzNotfallmedizin2}} 
%   \item[\citep{IntroClinEmergMed4}] \emph{\bibentry{IntroClinEmergMed4}} 
%   \item[\citep{Longmore2007}] \emph{\bibentry{Longmore2007}} 
%   \item[\citep{Fauci2008}] \emph{\bibentry{Fauci2008}}
% \end{labeling}
% }
% \end{multicols}
% \end{WidePar}